%---------------------------------------------------------------------------
% Cover-letter template matching the CV brand
%   - Same Plex stack (Serif body + Sans accent + Mono URLs)
%   - Same accent #0A2540
%   - Same fancyhdr footer with "Updated <date>" + page count
%   - Plain article class with sane margins for readable letter
%
% Customize the recipient + body each time you send. Keep it to ONE page
% per the v3 outreach research (Eugene Vinitsky, MIT UROP guidance).
%---------------------------------------------------------------------------
\documentclass[11pt,letterpaper]{article}

\usepackage{cmap}
\usepackage[T1]{fontenc}
\usepackage[utf8]{inputenc}
\usepackage[margin=1.0in,top=0.7in,bottom=0.85in]{geometry}

\usepackage[scale=0.97]{plex-serif}
\usepackage[scale=0.91]{plex-sans}
\usepackage[scale=0.89]{plex-mono}
\usepackage[protrusion=true,expansion=true,tracking=true,kerning=true,final]{microtype}
\AtBeginDocument{\figureversion{osf}}
\frenchspacing

\linespread{1.18}
\setlength{\parindent}{0pt}

\usepackage{xcolor}
\definecolor{accent}{HTML}{0A2540}
\definecolor{accentlight}{HTML}{2563A8}
\definecolor{body}{HTML}{1A1A1A}
\definecolor{meta}{HTML}{6B6B6B}
\definecolor{rule}{HTML}{C8C8C8}
\color{body}

\usepackage[
  pdfusetitle,
  pdfauthor={Dennis Johan Loevlie},
  pdftitle={Cover Letter --- Dennis Johan Loevlie},
  pdflang=en-US,
  hidelinks,
  colorlinks=true, urlcolor=accentlight, linkcolor=accentlight
]{hyperref}
\usepackage{hyperxmp}
\usepackage{xurl}
\usepackage{parskip}
\setlength{\parskip}{0.6\baselineskip plus 0.1\baselineskip}

\usepackage{fancyhdr}
\pagestyle{fancy}
\fancyhf{}
\renewcommand{\headrulewidth}{0pt}
\renewcommand{\footrulewidth}{0pt}
\fancyfoot[L]{\sffamily\scriptsize\color{meta}\today}%
\fancyfoot[R]{\sffamily\scriptsize\color{meta}%
  Dennis Loevlie \,\textperiodcentered\, %
  \href{https://dennisloevlie.com}{dennisloevlie.com} \,\textperiodcentered\, %
  CV: \href{https://loevlie.github.io/cv/loevlie-cv-latest.pdf}{loevlie.github.io/cv}}

\newcommand{\bdash}{\,\textemdash\,}

\begin{document}

%---------------------------------------------------------------------------
% Sender header — mirror of CV header for brand continuity
%---------------------------------------------------------------------------
{\sffamily\fontsize{18pt}{22pt}\selectfont\bfseries\color{accent} Dennis Johan Loevlie}\\[2pt]
{\sffamily\footnotesize\color{meta}%
\href{mailto:dennis.loevlie@tufts.edu}{dennis.loevlie@tufts.edu} \,\textperiodcentered\,
\href{https://dennisloevlie.com}{dennisloevlie.com} \,\textperiodcentered\,
\href{https://github.com/loevlie}{github.com/loevlie}}\\[2pt]
{\color{rule}\rule{\linewidth}{0.5pt}}

\vspace{1.2em}

%---------------------------------------------------------------------------
% Recipient
%---------------------------------------------------------------------------
\today
\vspace{1em}

{[Recipient Name]}\\
{[Lab / Group]}\\
{[Institution]}\\

\vspace{0.5em}

%---------------------------------------------------------------------------
% Body
%---------------------------------------------------------------------------
Dear [Dr.~Surname],

\vspace{0.4em}

[Opening — one or two sentences on the specific reason you're writing.
Reference the work, position, or program by name. Avoid "I hope this
finds you well" or "I have long admired your work" — the
v3 outreach research is unanimous that those signal templated mass mail.]

\vspace{0.4em}

[Middle paragraph(s) — your strongest credential paragraphs. Lead with
specifics: which paper, which result, which technique. For each claim,
have a numeric anchor and an artifact (paper, repo, model). For an
academic-applicant letter, this is where the ELLIS PhD framing,
Hulsebos and van~de~Meent collaboration, and your tabular-FM focus go.]

\vspace{0.4em}

[Closing — a specific, time-bounded ask. ``Could we have a 20-minute
call in the week of [date]?'' beats ``I'd love to chat sometime.''
Mention attached/linked CV. Sign off plainly.]

\vspace{0.4em}

Sincerely,

\vspace{2em}

Dennis Loevlie\\
{\sffamily\footnotesize\color{meta}%
ELLIS PhD candidate \bdash{} CWI (TRLLab) \& UvA (AMLab) \\
Multimodal foundation models for tabular data \bdash{} probabilistic deep learning}

\end{document}
